\Article{20}
Whenever the number of conditional equations is very large, determining the correlates \(A,B,C\) etc. by direct elimination may become so lengthy that the person carrying out the calculation may lack the patience for it. It may then often be convenient to obtain a complete adjustment by successive approximations, with the aid of the theorem in the preceding article. Let the conditional equations be divided into two or more classes, and first determine an adjustment by which the equations of the first class are satisfied, neglecting the rest. Next, treat the observations altered by this adjustment so that only the equations of the second class are taken into account. Generally speaking, applying the second system of adjustments will disturb agreement with the equations of the first class; therefore, if only two classes have been made, we shall return to the equations of the first class and determine a third system that satisfies them; then we shall subject the observations, corrected three times, to a fourth adjustment,
\SourceLeaf{084}
where only the equations of the second class are considered. Thus, by turns attending now to the first class and now to the second, we shall obtain adjustments that continually decrease; and if the division has been arranged skillfully, after a few iterations we shall reach stable values. If more than two classes have been made, the procedure will be similar: the individual classes will successively enter into the calculation, after the last class, the first will enter the calculation again, and so on. But it suffices here to have indicated this method, whose effectiveness will certainly depend greatly on its skillful application.

\Article{21}
It remains for us to complete the proof of the lemma assumed in Art. 8, though for the sake of clarity we shall use other notations better suited to this subject.

Let \(x^0,x',x'',x'''\) etc. be indeterminates, and suppose that from the equations
\[
\begin{aligned}
n^{00}x^0+n^{01}x'+n^{02}x''+n^{03}x'''+\text{etc.}&=X^0\\
n^{10}x^0+n^{11}x'+n^{12}x''+n^{13}x'''+\text{etc.}&=X'\\
n^{20}x^0+n^{21}x'+n^{22}x''+n^{23}x'''+\text{etc.}&=X''\\
n^{30}x^0+n^{31}x'+n^{32}x''+n^{33}x'''+\text{etc.}&=X'''\\
\text{etc.}&
\end{aligned}
\]
elimination yields the following:
\[
\begin{aligned}
N^{00}X^0+N^{01}X'+N^{02}X''+N^{03}X'''+\text{etc.}&=x^0\\
N^{10}X^0+N^{11}X'+N^{12}X''+N^{13}X'''+\text{etc.}&=x'\\
N^{20}X^0+N^{21}X'+N^{22}X''+N^{23}X'''+\text{etc.}&=x''\\
N^{30}X^0+N^{31}X'+N^{32}X''+N^{33}X'''+\text{etc.}&=x'''\\
\text{etc.}&
\end{aligned}
\]

Now, substituting in the first and second equations of the second system the values of the quantities \(X,X',X'',X'''\) etc. from the first system, we obtain
\[
\begin{aligned}
x^0={}&N^{00}(n^{00}x^0+n^{01}x'+n^{02}x''+n^{03}x'''+\text{etc.})\\
&+N^{01}(n^{10}x^0+n^{11}x'+n^{12}x''+n^{13}x'''+\text{etc.})\\
&+N^{02}(n^{20}x^0+n^{21}x'+n^{22}x''+n^{23}x'''+\text{etc.})\\
&+N^{03}(n^{30}x^0+n^{31}x'+n^{32}x''+n^{33}x'''+\text{etc.})\quad\text{etc.}
\end{aligned}
\]
\SourceLeaf{085}
and also
\[
\begin{aligned}
x'={}&N^{10}(n^{00}x^0+n^{01}x'+n^{02}x''+n^{03}x'''+\text{etc.})\\
&+N^{11}(n^{10}x^0+n^{11}x'+n^{12}x''+n^{13}x'''+\text{etc.})\\
&+N^{12}(n^{20}x^0+n^{21}x'+n^{22}x''+n^{23}x'''+\text{etc.})\\
&+N^{13}(n^{30}x^0+n^{31}x'+n^{32}x''+n^{33}x'''+\text{etc.})\quad\text{etc.}
\end{aligned}
\]

Since each of these equations must plainly be an identity, in both the former and the latter we may substitute any determinate values whatever for \(x^0,x',x'',x'''\) etc. In the former let us substitute
\[
x^0=N^{10},\quad x'=N^{11},\quad x''=N^{12},\quad x'''=N^{13}\quad\text{etc.}
\]
and in the latter
\[
x^0=N^{00},\quad x'=N^{01},\quad x''=N^{02},\quad x'''=N^{03}\quad\text{etc.}
\]
After this, subtraction gives
\[
\begin{aligned}
N^{10}-N^{01}={}&(N^{00}N^{11}-N^{10}N^{01})(n^{01}-n^{10})\\
&+(N^{00}N^{12}-N^{10}N^{02})(n^{02}-n^{20})\\
&+(N^{00}N^{13}-N^{10}N^{03})(n^{03}-n^{30})\\
&+\text{etc.}\\
&+(N^{01}N^{12}-N^{11}N^{02})(n^{12}-n^{21})\\
&+(N^{01}N^{13}-N^{11}N^{03})(n^{13}-n^{31})\\
&+\text{etc.}\\
&+(N^{02}N^{13}-N^{12}N^{03})(n^{23}-n^{32})\\
&+\text{etc. etc.}
\end{aligned}
\]
This equation can also be written as
\[
N^{10}-N^{01}=\Sigma(N^{0\alpha}N^{1\beta}-N^{1\alpha}N^{0\beta})(n^{\alpha\beta}-n^{\beta\alpha})
\]
where \(\alpha\beta\) denote all combinations of unequal indices.

Hence it follows that if
\[
\begin{aligned}
&n^{01}=n^{10},\quad n^{02}=n^{20},\quad n^{03}=n^{30}\\
&n^{12}=n^{21},\quad n^{13}=n^{31},\quad n^{23}=n^{32},\quad\text{etc.}
\end{aligned}
\]
or, more generally, if \(n^{\alpha\beta}=n^{\beta\alpha}\), then we shall also have
\[
N^{10}=N^{01}
\]
\SourceLeaf{086}
And since the order of the indeterminates in the given equations is arbitrary, under that supposition it is plainly true in general that
\[
N^{\alpha\beta}=N^{\beta\alpha}
\]
